\section{Discussion}
\label{sec:discussion}

\subsection{Relation Semantic Insufficiency as a Unified Root Cause}

Run~014's relation semantics audit revealed that all five reviewed candidates carry
at least one relation edge where the KG's class-level encoding is insufficient to
validate a substrate-specific chemical or biological claim.
This finding unifies previously separate observations---B-1/B-2's directionality
problem, C-1's chemistry validity failure, and C-2's compartmentalisation risk---under
a single root cause:

\medskip
\noindent\textbf{Root cause}: The KG encodes relation types at the \textit{class level}
rather than the \textit{instance level}.
A relation of type \texttt{produces} between a reaction class and a functional group
class (e.g., $r$\_Methylation $\to$ $fg$\_Piperidine) asserts that the reaction
\textit{type} can produce compounds bearing the functional group in general, but does
not validate whether the specific substrate in the path actually undergoes this
transformation.
Similarly, a relation \texttt{is\_substrate\_of} between a metabolite and an enzyme
encodes shared-enzyme membership, not substrate-specific causal transmission.

\medskip
This is not a failure of the drift filter---the filter correctly identifies structural
relation categories.
It is a limitation of the KG's underlying schema: without substrate-specificity
metadata on \texttt{produces} edges, and without compartment annotations on nodes,
the pipeline cannot distinguish mechanistically valid paths from structurally valid
but biologically incorrect ones.
Table~\ref{tab:relation_semantics} summarises the five semantic categories
and their inferential properties.

% Table 3: Relation Semantics Categories
% Source: paper_assets/relation_semantics_table.csv (representative rows)
% Full table: 22 rows; only representative rows shown here
\begin{table}[h]
\caption{%
  \textbf{Relation type classification by semantic category (Run~014, representative rows).}
  Five categories determine the maximum inferential claim a composed path
  can support.
  A path's inferential power is bounded by its weakest relation category.
  CR relations require substrate-specificity pre-check before any mechanistic
  interpretation is warranted.
  Full 22-row table: \texttt{paper\_assets/relation\_semantics\_table.csv}.
}
\label{tab:relation_semantics}
\small
\begin{tabular}{llllc}
\toprule
Relation & Category & Domain & Direction semantics & Safe for mech.\ inference \\
\midrule
\texttt{encodes}              & DM & bio  & causal (gene$\to$protein)     & \checkmark \\
\texttt{activates}            & DM & bio  & causal (regulator$\to$target)  & \checkmark \\
\texttt{inhibits}             & DM & bio, chem & causal (suppressor$\to$target) & \checkmark (track polarity) \\
\texttt{catalyzes}            & DM & bio  & causal (enzyme$\to$substrate) & conditional \\
\texttt{produces} (bio)       & DM & bio  & causal (enzyme$\to$product)   & \checkmark \\
\addlinespace
\texttt{produces} (rxn$\to$fg) & CR & chem & class-level generalisation   & \textbf{$\times$} (substrate check) \\
\addlinespace
\texttt{is\_substrate\_of}    & DN & bio  & ordering (metabolite$\to$enzyme; inverted) & $\times$ \\
\texttt{undergoes}            & DN & bio,chem & ordering (molecule$\to$rxn class) & $\times$ \\
\texttt{is\_precursor\_of}    & DN & chem & ordering (precursor$\to$product)  & $\times$ \\
\texttt{is\_product\_of}      & DN & chem & ordering (reversed; product$\leftarrow$precursor) & $\times$ \\
\addlinespace
\texttt{contains}             & SS & chem & structural composition & $\times$ \\
\texttt{is\_isomer\_of}       & SS & chem & symmetric              & $\times$ \\
\texttt{is\_reverse\_of}      & SS & chem & symmetric              & $\times$ \\
\texttt{interacts\_with}      & SS & bio  & undirected             & $\times$ \\
\addlinespace
\texttt{reacts\_with}         & CD & chem & context-dependent      & $\times$ \\
\bottomrule
\end{tabular}

\smallskip
\noindent\textbf{Legend:}
DM = directional-mechanistic;
DN = directional-but-non-causal;
SS = symmetric-structural;
CD = context-dependent;
CR = chemically-risky.
Blocked by Run~012 filter: \texttt{contains}, \texttt{is\_product\_of},
\texttt{is\_reverse\_of}, \texttt{is\_isomer\_of}.
CR relation (\texttt{produces} rxn$\to$fg) passes the filter but requires
pre-validation.
\end{table}


\subsection{The Alignment Operator as the Necessary Structural Primitive}

The align operator's mechanism---collapsing bridge entities to create a shared node
that BFS can traverse---is the necessary structural primitive for cross-domain discovery.
Without it, the composed graph has no edges crossing the domain boundary,
and compose cannot generate cross-domain candidates regardless of graph size.

This is a structural necessity (Claim~1), not a tuning choice.
The operator is analogous to Swanson's ABC model but formalised as a graph-algebraic
step: the bridge entity $B$ is the aligned node that connects $A$ (a bio entity)
to $C$ (a chem entity) via a path that could not exist in either domain alone.

\subsection{Quality Filtering: Generalisation Evidence and Residual Precision Gap}

The filter's generalisation across three independent domain pairs (Run~013) is
the strongest evidence against post-hoc overfitting.
A filter derived from 20 labeled candidates in one domain pair
(cancer signaling / metabolic chemistry) achieves identical qualitative outcomes
in two entirely different chemical systems (eicosanoid / natural products;
neurotransmitter / neuro-pharmacology) without parameter adjustment.

This suggests the four excluded relation types
(\texttt{contains}, \texttt{is\_product\_of}, \texttt{is\_reverse\_of},
\texttt{is\_isomer\_of}) represent a principled domain-general category:
structural-chemical description relations that are true but non-mechanistic.
The category appears stable across bio-chem domain pairs.

However, C-1 exposes a residual precision gap: the filter cannot distinguish
substrate-specific \texttt{produces} (bio enzyme $\to$ specific metabolite; safe)
from class-level \texttt{produces} (reaction class $\to$ functional group class; risky).
Both receive the same strong-relation credit.
Future filter versions should add a \texttt{produces\_class} flag to distinguish
these patterns, and should flag paths containing \texttt{is\_precursor\_of}
or \texttt{is\_substrate\_of} as carrying DN-category semantics.

\subsection{Bridge Dispersion: Design Implication for KG Construction}

The bridge dispersion finding (Claim~3) has a direct actionable implication for
knowledge graph curation targeting cross-domain discovery:
\textit{prefer bridge entities that are metabolic hubs or signaling intermediates
with many known interactors in both domains} over simply increasing the number
of bridge pairs.

Subset~B's five eicosanoids yield 8$\times$ more candidates per bridge than
Subset~A's NADH because each eicosanoid connects to multiple downstream synthesis
enzymes (bio) and natural product structural families (chem).
This amplification effect is non-linear: a metabolically diverse bridge class
generates combinations that a hub-low bridge cannot, regardless of how many
low-diversity bridge pairs are added.

\subsection{Top-20 Dominance by 2-Hop Candidates}

In all subsets and all scoring configurations, the top-20 candidates are exclusively
2-hop paths (mean score 0.78 across subsets).
The H4 result (Run~010) shows that the revised scorer improves the relative ranking
of deep candidates (309 promoted vs.\ 209 demoted), but does not change top-20
composition.
This is a systematic limitation: the absolute scoring advantage of well-supported
2-hop chains exceeds the relative gain from the quality-aware traceability revision.

Deep cross-domain candidates---the scientifically more interesting outputs---require
either a depth-promotion bonus, a minimum-depth tier, or a separate ranked list to
surface in practical use.
This is an unsolved problem in the current pipeline.
